\chapter{Token-Management und Kosten}
\label{chap:token_management}

% ============================================================
% AUTOR-NOTIZ STATT LERNZIELE
% ============================================================
\autornotiz{
2023 habe ich für einen E-Commerce-Kunden ein LLM-System gebaut. 50k Produktdokumente, 200k Queries/Tag.
Erster Versuch: Naive RAG mit OpenAI Embeddings + Pinecone. Hit Rate @ 5: 67\%. Latenz p99: 3.2s. Kosten: \$450/Tag.
Nach 3 Wochen Tuning: Recursive Chunking (512/64) + Hybrid Search (BM25 + Embedding) + Cross-Encoder Re-Ranker + pgvector.
Hit Rate @ 5: 94\%. Latenz p99: 800ms. Kosten: \$38/Tag.
Der Unterschied war nicht das Embedding-Modell. Es war Disziplin bei Chunking, Retrieval und Evaluation.
}

% ============================================================
% MOTIVATION
% ============================================================
\section{Motivation}

Die meisten AI-Engineering-Fehler sind keine Architektur-Fehler --- sie sind
Kosten- und Kontext-Fehler. Ein Entwickler packt 50 Dokumente in einen
Prompt, weil das Modell 1M Token Kontext hat --- und wundert sich über eine
Rechnung von 5~€ pro Query. Ein anderes Team nutzt GPT-4o für jede Anfrage,
obwohl GPT-4o-mini in 80\,\% der Fälle genauso gut ist.

Token-Management ist die grundlegendste Engineering-Disziplin im Umgang mit
LLMs. Während Software-Engineer:innen normalerweise in Request-Volumen oder
Datenvolumen denken, wird bei LLMs in \textbf{Tokens abgerechnet}. Ein Token
ist die kleinste Verrechnungseinheit --- kein Request, kein Byte, kein Wort.

\par



\par\mbox{}\par
\begin{realitycheckEnv}
Der grösste Kostenfehler: Jeden Request mit dem teuersten Modell zu bearbeiten.
GPT-4o-mini reicht für 80\,\% aller Aufgaben. Ein Semantic Router, der einfache
Queries automatisch an günstigere Modelle leitet, spart 60--80\,\% der Kosten.
\end{realitycheckEnv}

Die Konsequenz: Ohne Token-Management weisst du weder, was ein API-Call kostet,
noch, ob der Prompt ins Context-Limit passt. In der Praxis führen falsche
Token-Kalkulationen zu:
\begin{itemize}[leftmargin=*]
      \item \textbf{Kosten-Explosion}: Ein Batch-Job, der auf 10K Requests ausgelegt war,
            kostet plötzlich das Zehnfache, weil jeder Prompt 4.000 statt 400 Tokens hatte.
      \item \textbf{Cutoff-Fehler}: Requests werden stillschweigend abgeschnitten,
            weil der Prompt das Context-Limit überschreitet --- der Fehler erscheint erst in der Response.
      \item \textbf{Inkonsistente Ergebnisse}: Unterschiedliche Tokenizer (OpenAI vs. Anthropic)
            liefern unterschiedliche Resultate bei gleichem Text.
      \item \textbf{Insecure Context}: Kontext als Angriffsfläche (Poisoning, Injection, Leakage).
\end{itemize}

Token-Management ist keine optionale Optimierung, sondern eine
\textbf{Grundvoraussetzung} für jede produktive LLM-Integration.

% ============================================================
% GRUNDLAGEN
% ============================================================
\section{Grundlagen -- Tokenisierung und Kontext}

\subsection{Was ist ein Token?}

Ein Token ist die kleinste Einheit, die ein LLM verarbeiten kann.
Anders als natürliche Sprache, die in Wörtern oder Zeichen denkt,
arbeitet ein Transformer auf einer \textbf{vordefinierten Vokabelliste}
von Tokens. Ein Token kann sein:
\begin{itemize}[leftmargin=*]
      \item Ein ganzes Wort (z.\,B. \texttt{the} = 1 Token)
      \item Ein Wortteil (z.\,B. \texttt{un} + \texttt{believe} + \texttt{able})
      \item Ein einzelnes Zeichen bei unbekannten Wörtern
      \item Whitespace und Satzzeichen (jeweils eigene Tokens)
\end{itemize}

Die Umwandlung von Text in Tokens heisst \textbf{Tokenisierung}.
Der Tokenizer ist ein fester Bestandteil jedes Modells --- er wird
beim Training mitgelernt und ändert sich nicht mehr.

\subsection{Die drei Haupt-Tokenisierer}

\begin{table}[H]
\centering
\begin{tabularx}{\linewidth}{lXX}
\toprule
\textbf{Modell-Familie} & \textbf{Tokenizer} & \textbf{Algorithmus} \\
\midrule
GPT-4o / GPT-4 / o1 & TikToken & BPE (Byte-Pair-Encoding) \\
Claude 3.5 / Claude 4 & Claude Tokenizer & SentencePiece (BPE) \\
Gemini 2.0 / 2.5 & Google SentencePiece & SentencePiece (Unigram) \\
LLaMA 3.x & LLaMA Tokenizer & BPE (gguf-optimiert) \\
\midrule
\end{tabularx}
\caption{Tokenizer der wichtigsten Modell-Familien}
\label{tab:tokenizer-vergleich}
\end{table}

\subsubsection{BPE (Byte-Pair-Encoding)}

BPE ist der dominierende Tokenisierungs-Algorithmus. Er funktioniert
in drei Schritten:
\begin{enumerate}[leftmargin=*]
      \item \textbf{Zerlegung}: Jedes Zeichen wird einzeln als Token behandelt
      \item \textbf{Merging}: Die häufigsten Paare werden iterativ zu einem Token zusammengefasst
      \item \textbf{Vokabular}: Nach einer festgelegten Anzahl von Merges entsteht ein Vokabular (GPT-4o: ca. 100K Tokens)
\end{enumerate}

\subsubsection{SentencePiece}

SentencePiece ist ein \textbf{sprachunabhängiger} Tokenizer, der direkt
auf rohem Text arbeitet, ohne Whitespace-Pre-Tokenisierung.
Er wird von Anthropic und Google verwendet.

\subsection{Tokenisierung im Vergleich: Deutsch vs. Englisch}

Deutsch erzeugt typischerweise \textbf{1,4--2x mehr Tokens als Englisch}
für denselben Inhalt. Der Grund liegt in deutschen Komposita:

\begin{table}[H]
\centering
\begin{tabularx}{\linewidth}{lXXX}
\toprule
\textbf{Wort} & \textbf{GPT-4o (BPE)} & \textbf{Tokens} & \textbf{Bemerkung} \\
\midrule
the cat & [the] [cat] & 2 & Kurze Wörter \\
Kraftfahrzeughaftpflicht & [K] [raft] [fahr] [zeug] [haf] [t] [pfl] [icht] & 8 & Kompositum zerfällt \\
Haftpflichtversicherung & [Haf] [t] [pfl] [icht] [vers] [ich] [er] [ung] & 8 & Langes Wort \\
\midrule
\end{tabularx}
\caption{Token-Zerlegung Deutsch vs. Englisch (gpt-4o)}
\label{tab:token-deutsch}
\end{table}

\praxishinweis{ Deutsche Texte kosten 30--50\,\% mehr als englische durch längere
Tokenisierung. Wenn du Kosten optimieren willst, schreibe englische Prompts
und antworte in Englisch --- die Qualitätsunterschiede sind bei modernen
Modellen minimal. }

\subsection{Das Context Window}

Das Context Window (Kontextfenster) ist die maximale Menge an Tokens,
die ein Modell in einem Durchlauf verarbeiten kann. Es umfasst:

\begin{itemize}[leftmargin=*]
      \item \textbf{System-Prompt}: Verhaltensdefinition, Regeln, Persona
      \item \textbf{User-Prompt}: Die aktuelle Anfrage
      \item \textbf{Chat-History}: Vorherige Nachrichten im Gespräch
      \item \textbf{RAG-Kontext}: Dokumente, die zur Beantwortung relevant sind
      \item \textbf{Output}: Die Antwort des Modells (auch Teil des Fensters)
\end{itemize}

Die Summe aller Komponenten darf das Context Window nicht überschreiten.

\begin{table}[H]
\centering
\begin{tabularx}{\linewidth}{lX}
\toprule
\textbf{Modell} & \textbf{Kontextfenster} \\
\midrule
GPT-4o / GPT-4o-mini & 128K Tokens \\
Claude 4 Sonnet / Opus & 200K Tokens \\
Gemini 2.5 Flash / Pro & 1M Tokens \\
Llama 3 (alle Grössen) & 128K Tokens \\
Mistral Large & 128K Tokens \\
Qwen 2.5 & 128K Tokens \\
DeepSeek-V3 & 128K Tokens \\
\bottomrule
\end{tabularx}
\caption{Kontextfenster der wichtigsten Modelle}
\label{tab:kontextfenster}
\end{table}

\subsection{Input- vs. Output-Tokens}

APIs berechnen Input- und Output-Tokens getrennt:

\begin{itemize}

      \item \textbf{Input-Tokens} Was du an das Modell schickst (Prompt + History)
            --- inklusive aller Context-Tokens. Typisch 70--90\,\% der
            Gesamtkosten bei RAG-Systemen.
      \item \textbf{Output-Tokens} Was das Modell generiert. Output-Tokens sind
            3--10x teurer als Input-Tokens, da sie rechenintensiver sind
            (autoregressive Generation).
\end{itemize}

\begin{verbatim}
     Gesamtkosten = Input-Tokens * Input-Preis
                     + Output-Tokens * Output-Preis
\end{verbatim}

\begin{table}[H]
\centering
\begin{tabularx}{\linewidth}{lrr}
\toprule
\textbf{Modell} & \textbf{Input (pro Mio. Tokens)} & \textbf{Output} \\
\midrule
GPT-4o-mini & 0,15~\$ & 0,60~\$ \\
GPT-4o & 2,50~\$ & 10,00~\$ \\
Claude 4 Sonnet & 3,00~\$ & 15,00~\$ \\
Claude 4 Opus & 15,00~\$ & 75,00~\$ \\
Gemini 2.5 Flash & 0,08~\$ & 0,30~\$ \\
Gemini 2.5 Pro & 1,25~\$ & 10,00~\$ \\
DeepSeek-V3 & 0,27~\$ & 1,10~\$ \\
Llama 3 (lokal) & 0,00~\$ (Strom) & 0,00~\$ (Strom) \\
\bottomrule
\end{tabularx}
\caption{API-Kosten pro Million Tokens}
\label{tab:api-kosten}
\end{table}

\praxishinweis{ Output-Tokens sind 3--10x teurer als Input-Tokens. Setze \texttt{max\_tokens}
immer auf das Minimum, das deine Anwendung benotigt. Ein unlimitiertes
Output-Token kann bei langen Generierungen das Budget sprengen. }

\subsection{Lost in the Middle}

Das wichtigste Kontext-Phänomen: Modelle schenken dem Anfang und Ende
eines langen Prompts mehr Aufmerksamkeit als der Mitte. Studien zeigen,
dass die Retrieval-Genauigkeit in der Mitte eines 100K-Kontexts um
30--50\,\% niedriger sein kann als am Anfang oder Ende.

\begin{verbatim}
     Retrieval-Genauigkeit ueber die Kontext-Position:

      100% |\              /|
           | \            / |
       80%  |   \          /   |
           |    \        /    |
       60%  |     \------/     |
           | Anfang  Mitte  Ende
       40%  +-------------------+
               Kontext-Position
\end{verbatim}

Konsequenz: Auch bei grossen Kontextfenstern solltest du die wichtigsten
Informationen an den Anfang oder das Ende des Prompts setzen. Relevante
Dokumente gehoren in die ersten 20\,\% oder letzten 10\,\% des Kontexts.

% ============================================================
% ARCHITEKTUR
% ============================================================
\section{Architektur -- Context-Management in Produktion}

Jeder LLM-Call durchläuft ein Context-Budgeting:

\begin{verbatim}
     Request kommt an
           |
          v
     [Context-Budget: max 100K Tokens]
           |
     +----+----+
     |          |
    RAG      Chat-History
     (50K)     (30K)
     |          |
     +----+----+
          |
     System-Prompt (5K)
          |
     User-Input (1K)
          |
     Summe: 86K < 100K => OK
\end{verbatim}

\subsection{Context-Window-Management-Strategien}

\begin{itemize}

      \item \textbf{Sliding Window} Behalte nur die letzten N Nachrichten.
            Älteste Nachrichten werden verworfen. Einfach, effektiv, aber
            verliert Kontext bei langen Gesprächen.
      \item \textbf{Summary Compression} Fasse alte Nachrichten in einer
            Kurzzusammenfassung zusammen. Teurer (extra LLM-Call), aber
            erhält den Kontext.
      \item \textbf{Semantic Filtering} Behalte nur Nachrichten, die für die
            aktuelle Query relevant sind. Nutzt Embedding-Ähnlichkeit zur
            Selektion. Aufwändig, aber am besten für RAG-Systeme.
      \item \textbf{Hybrid} Sliding Window für kurze Konversationen + Summary
            Compression für lange Sessions. Der Standard.
\end{itemize}

% ============================================================
% CACHING UND KOSTEN-OPTIMIERUNG
% ============================================================
\section{Caching und Kosten-Optimierung}

\subsection{Prompt Caching (KV-Cache)}

Der effektivste Caching-Mechanismus ist der KV-Cache. Jeder LLM-Call
berechnet während der Prefill-Phase Key- und Value-Vektoren für alle
Input-Token. Wenn ein Teil des Inputs exakt gleich bleibt (z.\,B. der
System-Prompt), kann dieser Teil des KV-Cache wiederverwendet werden:

\begin{lstlisting}[language=Python, caption={Prompt Caching mit Anthropic und OpenAI}]
# Anthropic: Prompt Caching via cache_control
response = client.messages.create(
    model="claude-sonnet-4-20250514",
    max_tokens=1024,
    system=[
         {
             "type": "text",
             "text": SYSTEM_PROMPT,   # 2.000 Tokens, statisch
             "cache_control": {"type": "ephemeral"},
         }
     ],
    messages=[{"role": "user", "content": user_query}],
)
print(f"Cache: {response.usage.cache_creation_input_tokens > 0}")

# OpenAI: Prompt Caching (automatisch ab 1.024 Tokens Prefix)
response = client.chat.completions.create(
    model="gpt-4o",
    messages=[
         {"role": "system", "content": SYSTEM_PROMPT},
         {"role": "user", "content": user_query},
     ],
)
print(f"Cache Hit: {response.usage.prompt_tokens_details.cached_tokens > 0}")
\end{lstlisting}

Prompt Caching spart 50--90\,\% der Prefill-Zeit und bis zu 50\,\%
der Input-Kosten (Provider berechnen gecachte Tokens günstiger).

\subsection{Chat-History-Management}

Chat-History wächst unbegrenzt, wenn du sie nicht komprimierst.
Ab 80\,\% des Context-Limits sollte die History komprimiert werden:

\begin{itemize}[leftmargin=*]
      \item \textbf{Sliding Window}: N letzte Nachrichten behalten, alte verwerfen
      \item \textbf{Summarization}: Alte Nachrichten zusammenfassen (extra LLM-Call)
      \item \textbf{Hierarchisch}: Kurzform für frühe, Langform für spätere Nachrichten
\end{itemize}

% ============================================================
% SICHERHEIT
% ============================================================
\section{Sicherheit -- Kontext als Angriffsfläche}

\begin{itemize}[leftmargin=*]
      \item \textbf{Kontext-Injection ueber lange Prompts}: Ein Angreifer
            kann durch sehr lange Inputs das Context Window fullen und so
            die Kosten in die Hoehe treiben (Cost-Exploit). Setze harte
            Limits pro Request.
      \item \textbf{History-Poisoning}: Wenn User die Chat-History
            manipulieren können (z.\,B. durch vorherige Prompt Injection),
            kann das den Kontext vergiften. Validiere die History vor
            jedem neuen Call.
      \item \textbf{PII im Kontext}: Personenbezogene Daten, die im
            RAG-Kontext oder in der Chat-History enthalten sind, werden
            an die API gesendet. Maskiere PII vor dem Call.
      \item \textbf{Cache-Leakage}: Gecachte Prompts konnten sensible
            Informationen enthalten (Prompt-Engineering-Geheimnisse,
            System-Prompt mit Geschaeftslogik). Teile Caches nie zwischen
            verschiedenen Mandanten.
\end{itemize}

% ============================================================
% ZUSAMMENFASSUNG
% ============================================================
\begin{zusammenfassungEnv}
\begin{itemize}[leftmargin=*]
      \item Ein Token ist kein Wort --- deutsche Texte produzieren etwa 1,4--2x mehr Tokens als englische (Komposita!)
      \item Kontextfenster und Tokenisierung sind die Grundlage fur alle Kosten- und Architekturentscheidungen
      \item TikToken ist das Werkzeug zur Token-Zahlung: Ohne es tappst du im Dunkeln bei Kosten und Limits
      \item Output-Tokens sind 3--10x teurer als Input-Tokens --- priorisiere kurze Antworten
      \item Prompt-Caching (OpenAI + Anthropic) kann Input-Kosten um bis zu 90\,\% reduzieren
      \item Lost in the Middle: Wichtige Informationen gehoren an den Anfang oder das Ende des Kontexts
      \item Kontext als Angriffsfläche: Poisoning, Injection und Cache-Leakage sind real
      \item Grösseres Kontextfenster ist kein Ersatz fur gutes RAG. Vorauswahl ist besser als Volltext.
\end{itemize}
\end{zusammenfassungEnv}




\par\mbox{}\par
\begin{merkeEnv}
\begin{itemize}
      \item \textbf{Kontext ist teuer}: Jeder Token kostet Geld.
      \item \textbf{Output ist teurer als Input}: 3--10x Faktor.
      \item \textbf{Caching ist der effektivste Sparhebel}.
      \item \textbf{Grosser Kontext != gute Antwort}: Lost in the Middle.
      \item \textbf{Modell-Routing spart 80\,\% Kosten}.
\end{itemize}
\end{merkeEnv}

% ============================================================
% PRAXISPROJEKT -- KOSTENOPTIMIERUNGS-DASHBOARD
% ============================================================
\section{Praxisprojekt -- Kostenoptimierungs-Dashboard}

\textbf{Ziel:} Baue ein Dashboard, das die LLM-Kosten deiner Anwendung
in Echtzeit verfolgt und Optimierungsvorschlaege macht.

\textbf{Aufgaben:}
\begin{enumerate}[leftmargin=*]
      \item Implementiere einen CostTracker, der jeden API-Call mit
            Modell, Token-Zahlen und Kosten protokolliert
      \item Baue ein Semantic Router-System, das Queries automatisch
            an das gunstigste ausreichende Modell routet
      \item Erstelle einen Report, der zeigt: Kosten pro Tag, pro Modell,
            pro User, Kosten-Entwicklung ueber Zeit
      \item Fuge Budget-Alarme hinzu (80\,\%-, 90\,\%-, 100\,\%-Schwellen)
      \item Optimierte ein existierendes System um mindestens 50\,\%
            Kostenreduktion durch: kleineres Modell + Caching + Output-Limits
\end{enumerate}

\textbf{Testdaten:} Simuliere 1.000 API-Calls mit verschiedenen Modellen
und Prompt-Längen. Zeige die Kostenunterschiede der Optimierungsstufen.

\textbf{Erfolgskriterium:} Das Dashboard zeigt korrekte Tageskosten,
warnt bei Budget-Ueberschreitung und macht automatisch
Optimierungsvorschlaege.

% ============================================================
% WEITERFÜHRENDE RESSOURCEN
% ============================================================
\section{Weiterführende Ressourcen}

\begin{itemize}[leftmargin=*]
      \item \textbf{OpenAI Tokenizer}: \href{https://platform.openai.com/tokenizer}{platform.openai.com/tokenizer}
            --- Online-Tool zum Zählen und Analysieren von Tokens
      \item \textbf{tiktoken (GitHub)}: \href{https://github.com/openai/tiktoken}{github.com/openai/tiktoken}
            --- Pythons Tokenizer für OpenAI-Modelle
      \item \textbf{Anthropic Prompt Caching Guide}: \href{https://docs.anthropic.com/en/docs/build-with-claude/prompt-caching}{docs.anthropic.com}
      \item \textbf{OpenAI Pricing}: \href{https://openai.com/pricing}{openai.com/pricing}
      \item \textbf{Lost in the Middle Paper}: Liu et al. (2024):
            \glqq Lost in the Middle: How Language Models Use Long Contexts\grqq
      \item \textbf{Semantic Router (GitHub)}: \href{https://github.com/aurelio-labs/semantic-router}{github.com/aurelio-labs/semantic-router}
      \item \textbf{Langfuse Cost Tracking}: \href{https://langfuse.com/docs/cost-tracking}{langfuse.com/docs/cost-tracking}
\end{itemize}

\textbf{Nächster Schritt:} Jetzt kennst du die Token-Fakten. Im nächsten Kapitel geht es um Prompt Design -- wie du LLMs sprachlich steuern kannst.
\endinput

\clearpage
